% Generated draft: requires mathematical and source review.
\documentclass[11pt]{article}
\usepackage[margin=25mm]{geometry}
\usepackage{fontspec}
\setmainfont{Latin Modern Roman}
\usepackage{amsmath,amssymb,mathtools,unicode-math}
\setmathfont{Latin Modern Math}
\usepackage{enumitem,needspace,xurl,hyperref}
\setlength{\parindent}{0pt}
\setlength{\parskip}{5pt}
\setlength{\emergencystretch}{3em}
\newcommand{\N}{\mathbb{N}}
\newcommand{\Z}{\mathbb{Z}}
\newcommand{\Q}{\mathbb{Q}}
\newcommand{\R}{\mathbb{R}}
\newcommand{\C}{\mathbb{C}}
\newcommand{\F}{\mathbb{F}}
\newcommand{\A}{\mathbb{A}}
\newcommand{\PP}{\mathbb{P}}
\newcommand{\E}{\mathbb{E}}
\newcommand{\eps}{\varepsilon}
\newcommand{\dd}{\,\mathrm{d}}
\begin{document}
\section*{UnsolvedMath ID: 11000268}
\textbf{Source problem code:} catalogue entry
\subsection{Definitions and mathematical statement}
Let \(F_n\) be a finite-rank free group, \(n\geq2\), and let \(\phi\in\operatorname{Out}(F_n)\) be iwip (fully irreducible): no positive power of \(\phi\) preserves the conjugacy class of a proper nontrivial free factor. Outer Space \(CV_n\) in a standard projectivized model consists of marked finite metric graphs of rank \(n\), with positive edge lengths normalized to total length one, modulo marking-compatible isometry; equivalent graph conventions suppress valence-two vertices. The outer automorphism group acts by changing markings.

The source asks to describe the geometry of the axis bundle of \(\phi\), together with associated objects. It does not define that bundle or those objects. Common axis-related constructions involve folding lines associated with the attracting and repelling trees in the boundary of Outer Space, while displacement-minimizing loci depend on a chosen metric. These descriptions cannot silently be identified with the source's axis bundle, and no equality between them is asserted here.

A precise geometric answer must first specify the axis-bundle definition, projectivized versus unprojectivized model, topology and metric or folding-path structure, and then state the geometric/topological properties or invariants it describes. The Lipschitz geometry is asymmetric unless a symmetrization is chosen, so ``axis'' cannot be treated as an ordinary unique metric geodesic without further conditions. The source states iwip generally and does not restrict to atoroidal/ageometric cases; any such restriction needed for a particular construction would have to be made explicit. The rank range and desired completeness criterion for a description are otherwise absent.
\subsection{Short English statement}
Describe the axis-bundle geometry for a fully irreducible free-group outer automorphism. The precise bundle, metric and associated objects are not defined, so a particular folding-line or minimum-displacement model cannot be assumed.
\subsection{Sources}
\begin{itemize}
\item Local UnsolvedMath registry, permanent ID \(11000268\), exact source code \texttt{catalogue entry}.
\item Farb, \emph{Problems on Mapping Class Groups and Related Topics}, Question 1.3, PDF page 328; supplied running footer M. Bridson and K. Vogtmann: \url{https://www.math.uchicago.edu/~farb/papers/mcgbook.pdf}. Registry access 2026-07-29; pdf-font-structure extraction.
\item The summary refers generally to iwip/axis-bundle literature without any precise theorem or reference.
\end{itemize}
\end{document}
